\textbf{Phần 1. Trắc nghiệm lựa chọn}
\setcounterpageref{ex}{0}

\begin{ex} %[4D4V20]
Để cân khối lượng của một con voi trưởng thành, người ta thường dùng đơn vị nào?
\choice
{Gam}
{Mi-li-gam}
{\True Tấn}
{Yến}
\loigiai{
Con voi trưởng thành nặng vài tấn nên dùng đơn vị tấn.\\
\textbf{Đáp án của câu hỏi là C}
}
\end{ex}

\begin{ex} %[4D4V20]
Để đo diện tích của một phòng học tiểu học, đơn vị phù hợp nhất là
\choice
{$\text{cm}^2$}
{\True $\text{m}^2$}
{$\text{mm}^2$}
{$\text{km}^2$}
\loigiai{
Diện tích phòng học thường khoảng $40 - 60\text{ m}^2$.\\
\textbf{Đáp án của câu hỏi là B}
}
\end{ex}

\begin{ex} %[4D4V20]
Một bao xi măng nặng $50\text{ kg}$. $20$ bao xi măng như thế nặng bao nhiêu tấn?
\choice
{$0{,}5\text{ tấn}$}
{$2\text{ tấn}$}
{$10\text{ tấn}$}
{\True $1\text{ tấn}$}
\loigiai{
$20 \times 50 = 1000\text{ kg} = 1\text{ tấn}$.\\
\textbf{Đáp án của câu hỏi là D}
}
\end{ex}

\begin{ex} %[4D4V20]
Mỗi bước chân của bạn Nam dài khoảng $4\text{ dm}$. Nam đi từ cổng trường vào lớp hết $150$ bước. Quãng đường Nam đã đi là
\choice
{\True $60\text{ m}$}
{$600\text{ m}$}
{$6\text{ m}$}
{$6000\text{ m}$}
\loigiai{
$150 \times 4 = 600\text{ dm} = 60\text{ m}$.\\
\textbf{Đáp án của câu hỏi là A}
}
\end{ex}

\begin{ex} %[4D4V20]
Bạn Mai tập bơi cự li $50\text{ m}$ hết $48$ giây. Bạn Lan bơi hết $55$ giây. Bạn Mai bơi nhanh hơn bạn Lan bao nhiêu giây?
\choice
{$5$ giây}
{$8$ giây}
{\True $7$ giây}
{$10$ giây}
\loigiai{
$55 - 48 = 7$ giây.\\
\textbf{Đáp án của câu hỏi là C}
}
\end{ex}

\begin{ex} %[4D4V20]
Một mảnh đất trồng rau hình chữ nhật có chiều dài $12\text{ m}$, chiều rộng $5\text{ m}$. Diện tích mảnh đất đó là
\choice
{$34\text{ m}^2$}
{\True $60\text{ m}^2$}
{$600\text{ dm}^2$}
{$120\text{ m}^2$}
\loigiai{
$12 \times 5 = 60\text{ m}^2$.\\
\textbf{Đáp án của câu hỏi là B}
}
\end{ex}

\begin{ex} %[4D4V20]
Bác An thu hoạch lúa trên 2 thửa ruộng: Thửa 1 thu được $1$ tấn $2$ tạ, thửa 2 thu được $18$ tạ. Bác bán đi $2$ tấn thóc. Số thóc còn lại là
\choice
{$8$ tạ}
{$12$ tạ}
{$15$ tạ}
{\True $10$ tạ ($1$ tấn)}
\loigiai{
Tổng thu: $12 + 18 = 30$ tạ $= 3$ tấn.\\
Còn lại: $3 - 2 = 1$ tấn $= 10$ tạ.\\
\textbf{Đáp án của câu hỏi là D}
}
\end{ex}

\begin{ex} %[4D4V20]
Người ta quét vôi tường một phòng học có diện tích xung quanh là $80\text{ m}^2$. Tiền công quét vôi là $15\,000$ đồng mỗi mét vuông. Tổng số tiền công là
\choice
{\True $1\,200\,000$ đồng}
{$1\,000\,000$ đồng}
{$1\,500\,000$ đồng}
{$1\,800\,000$ đồng}
\loigiai{
$80 \times 15\,000 = 1\,200\,000$ đồng.\\
\textbf{Đáp án của câu hỏi là A}
}
\end{ex}

\begin{ex} %[4D4V20]
Một thang máy chở tối đa được $600\text{ kg}$. Có $8$ người lớn mỗi người nặng trung bình $70\text{ kg}$ cùng bước vào thang máy. Thang máy có bị quá tải không?
\choice
{Không bị quá tải}
{Vừa đủ tải trọng}
{\True Bị quá tải (nặng $560\text{ kg} < 600\text{ kg} \Rightarrow$ không quá tải, nếu 9 người thì quá tải)}
{Không xác định được}
\loigiai{
$8 \times 70 = 560\text{ kg} < 600\text{ kg}$ nên thang máy hoạt động bình thường an toàn.\\
\textbf{Đáp án của câu hỏi là C}
}
\end{ex}

\begin{ex} %[4D4V20]
Một tờ bìa hình chữ nhật có diện tích $6\text{ dm}^2$. Người ta cắt ra các hình vuông nhỏ cạnh $5\text{ cm}$. Cắt được tối đa bao nhiêu hình vuông nhỏ như thế?
\choice
{$20$ hình}
{\True $24$ hình}
{$25$ hình}
{$30$ hình}
\loigiai{
$6\text{ dm}^2 = 600\text{ cm}^2$.\\
Diện tích 1 hình vuông nhỏ: $5 \times 5 = 25\text{ cm}^2$.\\
Số hình: $600 : 25 = 24$ hình.\\
\textbf{Đáp án của câu hỏi là B}
}
\end{ex}

\begin{ex} %[4D4V20]
Một bộ phim hoạt hình bắt đầu chiếu lúc $19$ giờ $30$ phút và kết thúc lúc $20$ giờ $15$ phút. Thời lượng của bộ phim là
\choice
{$40$ phút}
{$50$ phút}
{$60$ phút}
{\True $45$ phút}
\loigiai{
$20$ giờ $15$ phút $- 19$ giờ $30$ phút $= 45$ phút.\\
\textbf{Đáp án của câu hỏi là D}
}
\end{ex}

\begin{ex} %[4D4V20]
Một xe tải có tải trọng cho phép chở là $3$ tấn hàng. Trên xe đã có sẵn $15$ tạ hàng. Người ta muốn xếp thêm các kiện hàng nặng $50\text{ kg}$ lên xe. Hỏi có thể xếp thêm tối đa bao nhiêu kiện hàng?
\choice
{\True $30$ kiện}
{$25$ kiện}
{$35$ kiện}
{$40$ kiện}
\loigiai{
Khối lượng còn chở thêm được: $30 - 15 = 15$ tạ $= 1500\text{ kg}$.\\
Số kiện hàng: $1500 : 50 = 30$ kiện.\\
\textbf{Đáp án của câu hỏi là A}
}
\end{ex}

\textbf{Phần 2. Trắc nghiệm đúng sai}
\setcounterpageref{ex}{0}

\begin{ex} %[4D4V20]
Xét tính hợp lí của các đơn vị đo trong thực tế:
\choiceTF
{\True Cân nặng một học sinh lớp 4 khoảng $35\text{ kg}$.}
{Diện tích sân trường tiểu học khoảng $1500\text{ cm}^2$.}
{\True Diện tích sân trường tiểu học khoảng $1500\text{ m}^2$.}
{\True Thời gian giải một bài toán khoảng $5$ phút.}
\loigiai{
\begin{itemize}
    \item Học sinh lớp 4 nặng ~35 kg: Đúng.
    \item Sân trường phải tính theo $\text{m}^2$ (không phải $\text{cm}^2$). Sai.
    \item Sân trường ~1500 $\text{m}^2$: Đúng.
    \item Giải 1 bài toán ~5 phút: Đúng.
\end{itemize}
\begin{itemize}
    \item a) \textbf{Đúng}.
    \item b) \textbf{Sai}.
    \item c) \textbf{Đúng}.
    \item d) \textbf{Đúng}.
\end{itemize}
}
\end{ex}

\begin{ex} %[4D4V20]
Một thửa ruộng thu hoạch được $2$ tấn $4$ tạ thóc. Người ta đã bán đi $\dfrac{1}{3}$ số thóc đó. Xét tính đúng sai:
\choiceTF
{\True Tổng số thóc thu hoạch được là $24$ tạ.}
{\True Số thóc đã bán là $8$ tạ.}
{Số thóc còn lại là $18$ tạ.}
{\True Số thóc còn lại là $16$ tạ ($1$ tấn $6$ tạ).}
\loigiai{
\begin{itemize}
    \item $2$ tấn $4$ tạ $= 24$ tạ. Đúng.
    \item Đã bán: $24 : 3 = 8$ tạ. Đúng.
    \item Còn lại: $24 - 8 = 16$ tạ $\neq 18$ tạ. Ý 3 Sai, Ý 4 Đúng.
\end{itemize}
\begin{itemize}
    \item a) \textbf{Đúng}.
    \item b) \textbf{Đúng}.
    \item c) \textbf{Sai}.
    \item d) \textbf{Đúng}.
\end{itemize}
}
\end{ex}

\textbf{Phần 3. Trả lời ngắn}
\setcounterpageref{ex}{0}

\begin{ex} %[4D4V20]
Một bao gạo nặng $5$ yến. Hỏi $6$ bao gạo như thế nặng bao nhiêu ki-lô-gam?
\shortans{300}
\loigiai{
\textbf{Đáp số:} $300$\\
$6 \times 50 = 300\text{ kg}$.
}
\end{ex}

\begin{ex} %[4D4V20]
Một mặt bàn hình chữ nhật dài $12\text{ dm}$, rộng $8\text{ dm}$. Tính diện tích mặt bàn theo đề-xi-mét vuông.
\shortans{96}
\loigiai{
\textbf{Đáp số:} $96$\\
$12 \times 8 = 96\text{ dm}^2$.
}
\end{ex}

\begin{ex} %[4D4V20]
Đổi $3$ thế kỉ $15$ năm ra năm.
\shortans{315}
\loigiai{
\textbf{Đáp số:} $315$\\
$315$ năm.
}
\end{ex}

\begin{ex} %[4D4V20]
Một người đi bộ với vận tốc $1\text{ m/s}$. Hỏi trong $5$ phút người đó đi được bao nhiêu mét?
\shortans{300}
\loigiai{
\textbf{Đáp số:} $300$\\
$5$ phút $= 300$ giây $\Rightarrow 1 \times 300 = 300\text{ m}$.
}
\end{ex}

\begin{ex} %[4D4V20]
Một bức tranh hình chữ nhật có diện tích $2\text{ m}^2$, chiều dài $20\text{ dm}$. Tính chiều rộng của bức tranh theo đề-xi-mét.
\shortans{10}
\loigiai{
\textbf{Đáp số:} $10$\\
$2\text{ m}^2 = 200\text{ dm}^2 \Rightarrow 200 : 20 = 10\text{ dm}$.
}
\end{ex}

\begin{ex} %[4D4V20]
Một xe tải có trọng tải $5$ tấn. Xe chở $30$ thùng hàng loại $1$ tạ/thùng và một số bao hàng loại $50\text{ kg}$/bao thì vừa đầy tải trọng. Tính số bao hàng loại $50\text{ kg}$.
\shortans{40}
\loigiai{
\textbf{Đáp số:} $40$\\
30 thùng $= 3000\text{ kg}$. Còn lại: $5000 - 3000 = 2000\text{ kg}$.\\
Số bao: $2000 : 50 = 40$ bao.
}
\end{ex}
